% --- 評価実験 ---
\section{評価実験}
\begin{itemize}
    \item 再現性を意識する．読者がこの章を読めばプログラムをゼロから実装できるレベルで詳しく書く．（ソースコードは論文に直貼りせず，GitHub等にアップロードしてリンクを掲載するほうがいいかも．）
    \item 自分の手法だけ載せるのではなく，既存手法と比較して優れていることを示す．
\end{itemize}

\subsection{実験}

\subsubsection{実験器具}
ここに実験環境などを書く．使用したデータセット，ハイパーパラメータ，計算機環境（CPU, GPU等），評価指標などを詳細に．

\subsubsection{実験手法}
ここに実験の手順を具体的に書く．

\subsection{結果}
結果を定量的に図や表にまとめる．表やグラフを見て分かる事実を述べる（例：「手法Aより手法Bの方が精度が5\%高かった」）．

\subsection{考察}
なぜそうなったのかを分析する（重要）．

「なぜ提案手法はうまくいった？」「逆にうまくいかなかったケースはある？それはなぜ？」など．これがないと実験レポート止まりになる．